\DSource{0130}{41}
\hypertarget{AB01-PDF0130-P01}{}
\DLine{AB01-PDF0130-body-L001}{}{}{diei quae 365 diebus addebatur. Qua re de motu Solis non dubitavit, ita ut Solem aliam}
\DLine{AB01-PDF0130-body-L002}{}{}{sphaeram habere cuius centrum a centris duarum sphaerarum differret, supponeret ({\itshape b}).}
\hypertarget{AB01-PDF0130-P02}{}
\DLine{AB01-PDF0130-body-L003}{}{}{\Indent Veteres plerumque haec ex observationibus aestivis, quae fiunt cum Sol per punctum}
\DLine{AB01-PDF0130-body-L004}{}{}{solstitiale aestivum transit, deprehenderunt; sed non ita verae videntur ut observationes factae}
\DLine{AB01-PDF0130-body-L005}{5}{}{cum Sol per alterutrum punctum aequinoctiale transit, et praesertim per punctum aequinoctii}
\DLine{AB01-PDF0130-body-L006}{}{p. 62.}{autumni; hoc enim tempore aer clarior et purior est quam tempore aequinoctii vernalis. Et}
\DLine{AB01-PDF0130-body-L007}{}{}{quia Solis motus in declinatione, quando per punctum solstitiale transit, tardus est, at quando}
\DLine{AB01-PDF0130-body-L008}{}{}{per puncta aequinoctialia iter facit, velocissimus, Ptolemaeus nonnisi observationibus autumnis}
\DLine{AB01-PDF0130-body-L009}{}{}{confisus est, et iuxta eas calculos instituit. Una ex observationibus Hipparchi qua usus est et}
\DLine{AB01-PDF0130-body-L010}{10}{}{de cuius veritate non dubitavit, ea est, qua, ut narrat (1), Solem per punctum aequinoctii}
\DLine{AB01-PDF0130-body-L011}{}{}{autumni transiisse comperit anno 178 ab Alexandro mortuo, die tertia ex quinque diebus epa-}
\DLine{AB01-PDF0130-body-L012}{}{}{gomenis, hora mediae noctis cuius crastinum fuit quarta epagomenorum dies, Alexandria in}
\DLine{AB01-PDF0130-body-L013}{}{}{urbe (2). Post hoc Ptolemaeus, 285 annis Aegyptiis transactis, iterum observavit, quam obser-}
\DLine{AB01-PDF0130-body-L014}{}{}{vationem ipse in libro suo accuratissimam fuisse dicit; et Solis transitum per punctum aequi-}
\DLine{AB01-PDF0130-body-L015}{15}{}{noctii autumni comperit anno tertio Antonini regis, id est anno 463 (3) ab Alexandro mor-}
\DLine{AB01-PDF0130-body-L016}{}{}{tuo, die nona mensis Coptici Athyr, una hora circiter post Solis ortum, Alexandria in urbe.}
\DLine{AB01-PDF0130-body-L017}{}{}{Intervallum inter duas observationes illas invenit esse veraciter 285 annorum Aegyptiorum,}
\DLine{AB01-PDF0130-body-L018}{}{}{70 dierum, et $\qfrac{1}{4}+\qfrac{1}{20}$ [$=\qfrac{3}{10}$] diei, pro 71 diebus $\qfrac{1}{4}$, qui ex quartis partibus integris in iis}
\DLine{AB01-PDF0130-body-L019}{}{}{285 annis colligi debebant. Ratio igitur huius unius diei minus $\qfrac{1}{20}$ diei, --- quo tempus ob-}
\DLine{AB01-PDF0130-body-L020}{20}{}{servationis praecessit tempus quartae partis diei 365 dies superantis, --- ad 285 annos inter}
\DLine{AB01-PDF0130-body-L021}{}{}{duas illas observationes transactos, est ut ratio unius diei ad 300 annos; qua re tempus unius}
\DLine{AB01-PDF0130-body-L022}{}{}{anni per has duas observationes inventum, complectitur 365 dies $\qfrac{1}{4}$, detracto $\qfrac{1}{300}$ diei, quod}
\DLine{AB01-PDF0130-body-L023}{}{}{est una pars et quinta e 360 partibus (4).}
\hypertarget{AB01-PDF0130-P03}{}
\DLine{AB01-PDF0130-body-L024}{}{}{\Indent Narrat etiam Ptolemaeus, se observationes antiquas aestivas, quae factae sint ante Hip-}
\DLine{AB01-PDF0130-body-L025}{25}{}{parchum, accepisse, e quibus observationem esse tempore Apseudis, regis [\textgreek{ἄρχων}] Athenarum,}
\DLine{AB01-PDF0130-body-L026}{}{}{factam, cum Sol per punctum solstitiale aestivum transivisset anno 108 (5) ante Alexandrum}
\DLine{AB01-PDF0130-body-L027}{}{}{mortuum, matutino XXI diei mensis Coptici Pharmouthi (6); deinde seipsum Solis transitum per}
\DLine{AB01-PDF0130-body-L028}{}{}{punctum solstitii aestivi observasse anno 463 (7) post Alexandrum mortuum, XI die mensis}
\DLine{AB01-PDF0130-body-L029}{}{}{Coptici Mesori, duabus horis circiter post mediam noctem cuius crastinum dies XII fuit (8).}
\DLine{AB01-PDF0130-body-L030}{30}{p. 63.}{Inter has duas observationes intercesserunt fere 571 anni Aegyptii, 140 dies et $\qfrac{1}{2}+\qfrac{1}{3}$}
\DLine{AB01-PDF0130-body-L031}{}{}{[$=\qfrac{5}{6}$] diei, pro 142 diebus $\qfrac{1}{2}+\qfrac{1}{4}$ [$=\qfrac{3}{4}$], qui ex quartis partibus in iis annis collecti}
\DLine{AB01-PDF0130-body-L032}{}{}{essent si in annis quartae partes integrae fuissent. Invenit igitur tempus solstitii aestivi tempus}
\DLine{AB01-PDF0130-body-L033}{}{}{quartae integrae partis diei uno die et $\qfrac{2}{3}+\qfrac{1}{4}$ [$=\qfrac{11}{12}$] praecessisse, quorum ratio ad 571}
\DLine{AB01-PDF0130-body-L034}{}{}{annos memoratos est ut ratio duorum integrorum dierum ad 600 (9) annos; idque cum eo}
\DLine{AB01-PDF0130-body-L035}{35}{}{iuxta quod calculos fecit convenit, scilicet singulis 300 annis tempus observationis tempus}
\DLine{AB01-PDF0130-body-L036}{}{}{quartae partis integrae diei uno die praecedere, licet ob causam memoratam hae observationes}
\DLine{AB01-PDF0130-body-L037}{}{}{aestivae non adeo sint bonae ut autumnae. Manifestum est inter primam illam observationem}
\par\vspace{6pt}\hrule height.25pt\vspace{5pt}
\noindent\begin{minipage}[t]{.48\textwidth}\vspace{0pt}\fontsize{9}{11.5}\selectfont\setlength{\GutterOffset}{0pt}
\hypertarget{AB01-PDF0130-N01}{}
\DLine{AB01-PDF0130-notes_left-L001}{}{}{\Indent (1) {\itshape Almag.} III, 2 (ed. Halma, t. I, p. 161).}
\hypertarget{AB01-PDF0130-N02}{}
\DLine{AB01-PDF0130-notes_left-L002}{}{}{\Indent (2) Al-Battānī coniectura; minime Ptolemaeus}
\DLine{AB01-PDF0130-notes_left-L003}{}{}{dicit Hipparchum id Alexandriae observasse. Et re}
\DLine{AB01-PDF0130-notes_left-L004}{}{}{vera nunc certum habetur Hipparchum Rhodi tan-}
\DLine{AB01-PDF0130-notes_left-L005}{}{}{tum observationes suas perfecisse, quam perperam}
\DLine{AB01-PDF0130-notes_left-L006}{}{}{sub eodem meridiano quo Alexandria iacere putabat.}
\hypertarget{AB01-PDF0130-N03}{}
\DLine{AB01-PDF0130-notes_left-L007}{}{}{\Indent (3) Codex perperam 563.}
\hypertarget{AB01-PDF0130-N04}{}
\DLine{AB01-PDF0130-notes_left-L008}{}{}{\Indent (4) $\dfrac{1+1/5}{360}=\dfrac{6/5}{360}=\dfrac{6}{1800}=\dfrac{1}{300}$.}
\end{minipage}
\hfill\begin{minipage}[t]{.48\textwidth}\vspace{0pt}\fontsize{9}{11.5}\selectfont\setlength{\GutterOffset}{0pt}
\hypertarget{AB01-PDF0130-N05}{}
\DLine{AB01-PDF0130-notes_right-L001}{}{}{\Indent (5) Codex perperam 168.}
\hypertarget{AB01-PDF0130-N06}{}
\DLine{AB01-PDF0130-notes_right-L002}{}{}{\Indent (6) Est observatio Metonis et Euctemonis (27}
\DLine{AB01-PDF0130-notes_right-L003}{}{}{Iun. 432 ante Chr. n.); {\itshape Almag.} III, 2 (ed. Halma,}
\DLine{AB01-PDF0130-notes_right-L004}{}{}{t. I, p. 162).}
\hypertarget{AB01-PDF0130-N07}{}
\DLine{AB01-PDF0130-notes_right-L005}{}{}{\Indent (7) Male Plato 36.}
\hypertarget{AB01-PDF0130-N08}{}
\DLine{AB01-PDF0130-notes_right-L006}{}{}{\Indent (8) {\itshape Almag.} l. l.}
\hypertarget{AB01-PDF0130-N09}{}
\DLine{AB01-PDF0130-notes_right-L007}{}{}{\Indent (9) Plato perperam 60.}
\end{minipage}
\par\vspace{5mm}\hfill{\fontsize{8}{10}\selectfont 6}\hspace{10mm}
